\section{Data}
\label{sec:data}

\subsection{Advertising Spending}

I have collected data on government advertising spending (COMSOC), which includes the date (at the daily level) each contract was enforced, the name of the media outlet receiving funds, the name of the government agency contracting the media services or purchasing airtime, and the amount of money allocated. Importantly, I will be able to locate the media outlets that produce newspaper articles using the COMSOC data base as each media outlet that received money is classified into radio, online, newspapers, magazines, billboards, etc. 

\subsection{Press conferences}

I have collected data on the official transcript of each morning press conference held by President López Obrador (a total of 1,423 conferences during his term). Additionally, I submitted an information request to the Office of the Presidency to obtain the names and media affiliations of all journalists attending each conference. This allows me to compile a comprehensive list of media outlets that have entered and participated in these press conferences. For initial analysis I have already extracted the information of journalists and media outlets that participated in the conference, using GPT and the structure of the conference. 

Moreover, I requested the results of the daily lottery implemented since April 2022, conducted before each morning press conference to assign journalists to the first, second, and third rows. Prior to 2022, seating was arranged on a first-come, first-served basis at the Palacio Nacional, which led to criticism as many journalists felt that the process was biased, favoring those aligned with the government over independent or opposition voices\citep{Tombola-Razon}. These journalists would often arrive as early as 1 a.m. to secure front-row seats. In response, the President proposed a lottery system to determine seating. By April 2022, the government formalized this change, announcing that journalists would be randomly assigned to the first three rows, with one journalist per row from each type of media—radio, television, print, etc \citep{Tombola-Financiero}.\footnote{Although this data currently appears to be restricted, reportedly for the safety of journalists, I have requested a censored version that excludes journalists' names, leaving only their media outlet associations, which is hardly sensitive information. I am actively working to contact the person responsible for the Presidency Office through connections I have within the Mexican federal government. If these efforts prove unsuccessful, I plan to reach out to Professor Luis Estrada from ITAM, who has been systematically collecting data from the conferences for his consulting firm, \href{https://spintcp.com/quienes-somos/}{SPIN}.}

\subsection{Newspaper data}

To compile a comprehensive dataset of newspaper articles that can be matched to this spending data, I will try to scrape a near-complete set of articles published online by newspapers from 2020 to 2024. I have started using web crawlers for data collection. At the moment I have fully scrapped the articles of 6 newspapers: La Jornada (top advertisement earner after the spending switch, left-leaning), Información Integral 24/7 (covers 5 states), AltaVoz (independent newspaper), Grupo Cantón (diario Basta, Quintana Roo Hoy) and Diario La Voz Oaxaca (local newspaper). Major media outlets like Reforma and Milenio have introduced paywalls on their websites, and some have complex HTML structures that make it difficult to extract articles directly. Therefore, for these more complex websites, following \cite{SobrinoCartel, Marshall} I will prompt Google News Advanced Search daily for the relevant time period and for each media outlet. I'll use common Spanish words in each prompt to recover all the articles published that particular day and then collect each URL from the several result pages. 

Ideally, from this dataset, I aim to calculate the proportion of critical and favorable government articles written by each media outlet $i$ during each time period $t$. Accurate labeling of the data is crucial for this purpose, and while there are several methods to approach this task, no single method is likely to be overwhelmingly superior to the others. Therefore, the initial analysis could use one method (the best one), with robustness checks performed using alternative approaches. Although any labeling method may introduce some degree of measurement error, as long as this error is not systematically correlated with the instrument or the dependent variable and it captures the variation across outlets and time effectively, the results would probably remain valid. 

Additionally, since the final variable used in the analysis will be aggregated across outlets and time periods, even imperfect measurements can still provide robust estimates if the chosen method effectively captures the overall variation in the desired variable. However, it is important to consider potential measurement challenges that could affect the interpretation of results. For instance, let us assume I aim to label violence-related articles as negative or critical, based on the rationale that violence is a highly salient and electorally costly topic. In this context, it would seem logical to assume that media outlets covering violence-related issues are inherently critical of the government. However, this assumption could lead to inaccuracies.

Consider an outlet aligned with the government, aware of the salience of this issue, which deliberately covers it positively. Such an outlet might publish numerous articles highlighting government successes—such as successful investigations, reductions in crime rates, or new public policies to combat violence—presenting the government in a favorable light. Meanwhile, a critical outlet might publish sensationalist stories daily, emphasizing violent incidents from across the country, even those unrelated to drug violence, to paint the situation as dire.

If both types of outlets use similar language—mentioning terms like violence, drugs, and deaths—a dictionary-based approach to classification might incorrectly interpret both as equally critical of the government. This would remove the true differences in tone and intent between the two outlets, misrepresenting their variation in reporting styles. Therefore, addressing this challenge requires careful consideration of the context and tone of articles, potentially through top of the line machine learning models or manual validation steps to ensure the labels accurately reflect the underlying intent of the reporting.

\subsubsection{BERT or Large Language Models (LLMs)}

For this task, the initial labeling process could leverage the metadata provided by many newspapers, where articles are often categorized into topics such as politics, economics, society, or violence. Some outlets even provide direct labels for security or violence-related articles. A practical first step would involve selecting a random sample from this labeled dataset to train a predictive model capable of classifying the remaining unlabeled articles.

For example, from La Jornada, I have already collected over 100,000 labeled articles. This dataset could serve as a starting point for training a model such as \href{https://github.com/dccuchile/beto}{BETO}, a Spanish-language BERT model that is straightforward to implement and has demonstrated strong prediction performance. Once a labeled dataset is established, and I have reduced the amount of articles to a relevant subset, more advanced approaches could be employed to refine predictions. For instance, fine-tuning a Large Language Model (LLM) specifically on topics such as violence, where context and tone are particularly important, could improve the classification of difficult cases. Proprietary models such as OpenAI's GPT, could be utilized depending on computational and financial resources.

\subsubsection{Rule-Based and Hybrid Models}

Going back to the example of outlets posting positively and negatively on violence, one solution to this could be adopting a hybrid approach combining rule-based and machine learning methods. For example, rules could be applied to pre-filter articles, reducing the noise in the dataset before training a BERT or LLM model. This two-step process can improve the model's accuracy, particularly when dealing with datasets where some articles are mislabeled or ambiguous.

Given this, rule-based methods could complement the analysis. For example, specific keywords or phrases commonly associated with criticism (e.g., "corruption," "failure," "mismanagement") or support (e.g., "achievement," "progress," "success") could be used to identify relevant articles. Or even restricting into salient and electorally damaging topics, such as violence or inflation. Some efforts have been made in the Mexican context using a dictionary based approach \citep{SobrinoCartel, Marshall} to identify articles about violence and cartel-presence.

To ensure robustness, different approaches to labeling—such as fine-tuning models on subsets of the data, applying semi-supervised learning, or testing across various pretrained models—should be evaluated. Cross-validation techniques could assess the consistency of the classifications across multiple subsets of the data. Additionally, testing the aggregation of critical and favorable proportions across outlets and time will help confirm that the method accurately captures temporal and cross-sectional variations.

By combining state-of-the-art models like BETO or LLMs with careful validation and supplementary rule-based approaches, this methodology will ensure the creation of a reliable variable, that could even be used for other unrelated projects.

